\section{简单闭曲线的环绕数}\label{ch10:section-winding-number}

若 $h\colon S^{1}\to\mathbb{R}^{2}-0$ 是一个连续映射, 则诱导出来的同态 $h_{*}$ 将 $S^{1}$ 的基本群的生成元映射为 $\mathbb{R}^{2}-0$ 的基本群的生成元的某整数幂. 这个整数幂 $n$ 就称为 $h$ 相对于 $0$ 的环绕数 (winding number). 它刻画了 $h$``绕着原点缠绕了 $S^{1}$''多少圈. 它的正负号当然依赖于生成元的选取. 参见 \cref{ch10:fig-65-1}. 在下一节中我们将对此予以更正式的介绍.

\par
\begin{figure}[htbp]
\centering
\includegraphics[width=0.7\textwidth]{ch10/fig-65-1.png}
\caption{环绕数 $n=\pm 2$ 与 $n=0$.}\label{ch10:fig-65-1}
\end{figure}

现在, 我们提出这样一个问题: 如果 $h$ 是一个单射, 也就是说, 如果 $h$ 是 $S^{1}$ 与 $\mathbb{R}^{2}-0$ 中一条简单闭曲线 $C$ 的同胚, 关于 $h$ 的环绕数会有什么结论呢? 根据 \cref{ch10:fig-65-2} 可以做出这样一个明显的猜测: 如果 $0$ 属于 $\mathbb{R}^{2}-C$ 的一个无界分支, 那么 $n=0$. 然而, 如果 $0$ 属于一个有界分支, 那么 $n=\pm 1$.

\par
\begin{figure}[htbp]
\centering
\includegraphics[width=0.7\textwidth]{ch10/fig-65-2.png}
\caption{单射情形下的猜测: $n=0$ 或 $n=\pm 1$.}\label{ch10:fig-65-2}
\end{figure}

第一个猜测是容易证明的, 因为根据零伦引理, 如果 $0$ 属于 $\mathbb{R}^{2}-C$ 的一个无界分支, 那么 $h$ 是零伦的. 另一方面, 第二个猜测却是惊人的困难, 实际上这是一个相当深刻的结论. 在本节中我们将证明它.

如常, 我们用 $S^{2}$ 代替 $\mathbb{R}^{2}\cup\{\infty\}$, 设 $p$ 为对应 $0$ 的点, $q$ 是对应 $\infty$ 的点. 这时我们的猜测可以重新整理为: 如果 $C$ 是 $S^{2}$ 中的一条简单闭曲线, 设 $p$ 和 $q$ 属于 $S^{2}-C$ 的不同分支, 那么内射 $j\colon C\to S^{2}-p-q$ 便诱导出基本群之间的一个同构. 这便是我们将要证明的结论.

首先, 我们证明上述结论在简单闭曲线 $C$ 包含于一个四顶点的完全图时成立. 然后再证明一般情况.

\begin{lemma}\label{ch10:lemma-winding-k4}
设 $G$ 是 $S^{2}$ 的一个子空间, 并且 $G$ 是以 $a_{1},\cdots,a_{4}$ 为顶点的完全图. 设 $C$ 是子图 $a_{1}a_{2}a_{3}a_{4}a_{1}$, 它是一条简单闭曲线. 设 $p$ 和 $q$ 分别表示两个边 $a_{1}a_{3}$ 和 $a_{2}a_{4}$ 的内点. 那么:

\begin{enumerate}[(a)]
\item 点 $p$ 和 $q$ 属于 $S^{2}-C$ 的不同分支.
\item 内射 $j\colon C\to S^{2}-p-q$ 诱导出基本群之间的一个同构.
\end{enumerate}
\end{lemma}

\begin{proof}
(a) 如引理 64.3 的证明中一样, $\theta$ 空间 $C\cup a_{1}a_{3}$ 将 $S^{2}$ 分割为三个分支 $U,V$ 和 $W$. 这三者之一, 譬如说 $W$, 将 $C$ 作为其边界. 这是唯一一条边界包含 $a_{2}$ 和 $a_{4}$ 的分支. 因此, $a_{2}a_{4}-a_{2}-a_{4}$ 必包含于 $W$, 于是 $q$ 属于 $W$. 当然 $p$ 不能属于 $W$, 因为 $p$ 属于 $\theta$ 空间 $C\cup a_{1}a_{3}$. 根据 $\theta$ 空间引理, $W$ 是 $S^{2}-C$ 的一个分支. 因此 $p$ 和 $q$ 属于 $S^{2}-C$ 的不同分支.

(b) 设 $X=S^{2}-p-q$. 证明思想如下: 选取点 $x$ 为弧 $a_{1}a_{2}$ 内部的一个点, 点 $y$ 为弧 $a_{3}a_{4}$ 内部的一个点, 令 $\alpha$ 和 $\beta$ 是折线路径
\[
\alpha=xa_{1}a_{4}y\quad\text{和}\quad\beta=ya_{3}a_{2}x,
\]
那么 $\alpha*\beta$ 是包含在简单闭曲线 $C$ 中的一条回路. 我们将要证明 $\alpha*\beta$ 代表 $X$ 的基本群的一个生成元. 由此得到同态 $j_{*}\colon\pi_{1}(C,x)\to\pi_{1}(X,x)$ 是满射, 从而 $j_{*}$ 是一个同构 (因为所涉及的群是无限循环群). 参见 \cref{ch10:fig-65-3}.

设 $D_{1}$ 和 $D_{2}$ 是弧
\[
D_{1}=pa_{3}a_{2}q\quad\text{和}\quad D_{2}=qa_{4}a_{1}p,
\]
令 $U=S^{2}-D_{1}$, $V=S^{2}-D_{2}$. 参见 \cref{ch10:fig-65-4}. 那么 $X=U\cup V$, 并且 $U\cap V=S^{2}-D$, 其中 $D$ 是简单闭曲线 $D=D_{1}\cup D_{2}$. 因此, 根据 Jordan 曲线定理, $U\cap V$ 有两个分支. 进而, 由于 $D$ 就是简单闭曲线 $a_{1}a_{3}a_{2}a_{4}a_{1}$, 所以根据结论 (a) 可见, 分别属于图 $G$ 另外两条边的内部的点 $x$ 和 $y$ 属于 $S^{2}-D$ 的两个不同分支.
\end{proof}

\par
\begin{figure}[htbp]
\centering
\includegraphics[width=0.6\textwidth]{ch10/fig-65-3-4.png}
\caption{左: 回路 $\alpha*\beta$; 右: 弧 $D_{1},D_{2}$ 与空间 $U,V$.}\label{ch10:fig-65-3-4}
\end{figure}

\begin{proof}[continued]
因此满足定理 63.1 的假设. $\alpha$ 是 $U$ 中从 $x$ 到 $y$ 的一条道路, $\beta$ 是 $V$ 中从 $y$ 到 $x$ 的一条道路. 因为 $X$ 的基本群是无限循环群, 所以回路 $\alpha*\beta$ 代表这个群中的一个生成元.
\end{proof}

\begin{theorem}\label{ch10:thm-winding-isomorphism}
设 $C$ 是 $S^{2}$ 中的一条简单闭曲线, $p$ 和 $q$ 属于 $S^{2}-C$ 的不同分支, 那么内射 $j\colon C\to S^{2}-p-q$ 诱导基本群之间的同构.
\end{theorem}

\begin{proof}
证明过程中需要构造一个以 $C$ 为子图的四顶点完全图.

\textbf{第一步.} 设 $a,b$ 和 $c$ 是 $\mathbb{R}^{2}$ 中的三个互异点. 如果 $A$ 是以 $a$ 和 $b$ 为端点的弧, $B$ 是以 $b$ 和 $c$ 为端点的弧, 那么便有一条包含在 $A\cup B$ 中的弧以 $a$ 和 $c$ 为端点.

选取一条从 $a$ 到 $b$ 的道路 $f\colon I\to A$ 和一条从 $b$ 到 $c$ 的道路 $g\colon I\to B$, 使得 $f$ 和 $g$ 都是同胚. 设 $t_{0}$ 是 $I$ 中满足条件 $f(t_{0})\in B$ 的最小点, $t_{1}$ 是 $I$ 中满足条件 $g(t_{1})=f(t_{0})$ 的点. 那么集合 $f([0,t_{0}])\cup g([t_{1},1])$ 就是所求的弧. (若 $t_{0}=0$ 或 $t_{1}=1$, 则这些集合中的某一个为单点集.) 参见 \cref{ch10:fig-65-5}.

\textbf{第二步.} 我们证明如果 $U$ 是 $\mathbb{R}^{2}$ 中的开集, 那么对于 $U$ 中的任何两个点, 如果它们能够由 $U$ 中的道路连接的话, 这两个点一定是 $U$ 中某一个弧的端点.

对于 $x,y\in U$, 定义 $x\sim y$, 如果 $x=y$ 或存在 $U$ 中的道路以 $x$ 和 $y$ 为端点. 第一步的结论说明这是一个等价关系. 等价类是开集, 因为如果 $x$ 的 $\varepsilon$-邻域包含在 $U$ 中, 那么这个邻域就是由等价于 $x$ 的点构成的. 因为 $U$ 是连通的, 故只有一个这样的等价类.

\textbf{第三步.} 设 $C$ 是 $\mathbb{R}^{2}$ 中的一条简单闭曲线. 我们构造 $\mathbb{R}^{2}$ 中的一个子空间 $G$, 它是以 $a_{1},\cdots,a_{4}$ 为四个顶点的完全图, 并且 $C$ 等于子图 $a_{1}a_{2}a_{3}a_{4}a_{1}$.

为方便起见, 我们假设 $0$ 属于 $\mathbb{R}^{2}-C$ 的一个有界分支. 考虑 $\mathbb{R}^{2}$ 中的 $x$ 轴 $\mathbb{R}\times 0$. 令 $a_{1}$ 表示 $x$ 轴的负半轴中包含在 $C$ 中的最大点, $a_{3}$ 表示 $x$ 轴的正半轴中包含在 $C$ 中的最小点, 那么线段 $a_{1}a_{3}$ 包含在 $\mathbb{R}^{2}-C$ 的有界分支的闭包中.

将 $C$ 表示为两个端点为 $a_{1}$ 和 $a_{3}$ 的弧 $C_{1}$ 和 $C_{2}$ 的并. 设 $a$ 是 $\mathbb{R}^{2}-C$ 的无界分支的一个点. 因为 $C_{1}$ 和 $C_{2}$ 不能将 $\mathbb{R}^{2}$ 分割, 所以我们能选取从 $a$ 到 $0$ 的两条道路 $\alpha\colon I\to\mathbb{R}^{2}-C_{1}$ 和 $\beta\colon I\to\mathbb{R}^{2}-C_{2}$, 根据第二步, 我们可以假设 $\alpha$ 和 $\beta$ 是单射. 设 $a_{2}=\alpha(t_{0})$, 其中 $t_{0}$ 是满足 $\alpha(t_{0})\in C$ 的最小的数, 那么 $a_{2}$ 是 $C_{2}$ 的一个内点. 类似地, 设 $a_{4}=\beta(t_{1})$, 其中 $t_{1}$ 是满足 $\beta(t_{1})\in C$ 的最小点, 那么 $a_{4}$ 是 $C_{1}$ 的内点. 于是 $\alpha([0,t_{0}])$ 和 $\beta([0,t_{1}])$ 分别是连接 $a$ 与 $a_{2}$ 和 $a_{4}$ 的弧. 由第二步, 它们的并包含一段端点为 $a_{2}$ 和 $a_{4}$ 的弧, 并且这段弧与 $C$ 仅交于这两个点. 这段弧以及线段 $a_{1}a_{3}$ 和曲线 $C$ 一起构成了所求的图. 参见 \cref{ch10:fig-65-6}.

\textbf{第四步.} 根据第三步的结论和前述引理可得: 对于包含在 $S^{2}-C$ 的不同分支中的某两个点 $p$ 和 $q$, 内射 $j\colon C\to S^{2}-p-q$ 诱导出基本群之间的同构. 为完成定理的证明, 我们只需要证明的是: 对于包含在 $S^{2}-C$ 的不同分支中的任意两个点 $p$ 和 $q$ 结论都成立.

为此, 需要证明以下事实:

设 $D$ 是 $\mathbb{R}^{2}$ 中的简单闭曲线, $0$ 属于 $\mathbb{R}^{2}-D$ 的一个有界分支. $p$ 是这个分支中的另一个点. 如果内射 $j\colon D\to\mathbb{R}^{2}-0$ 诱导出基本群之间的同构, 那么内射 $k\colon D\to\mathbb{R}^{2}-p$ 也诱导出基本群之间的同构.
\end{proof}

\par
\begin{figure}[htbp]
\centering
\includegraphics[width=0.8\textwidth]{ch10/fig-65-5-6.png}
\caption{左: 第一步的构造; 右: 以 $C$ 为子图的完全图 $G$.}\label{ch10:fig-65-5-6}
\end{figure}